% !TEX program = lualatex
% !TeX encoding = UTF-8
\documentclass[12pt,a4paper]{article}

% --- LuaLaTeX: 中文支持通过 ctex ---
\usepackage[fontset=none]{ctex}
\setmainfont{Liberation Serif}
\setCJKmainfont{SimSun}
\setCJKmonofont{SimSun}

% --- 数学和格式包 ---
\usepackage{amsmath,amssymb,amsthm}
\usepackage{graphicx}
\usepackage{hyperref}
\usepackage{url}
\usepackage{xcolor}
\usepackage{booktabs}
\usepackage{geometry}
\geometry{margin=2.5cm}

% ============================================================
% 自定义 YUCT 命令
% ============================================================
\newcommand{\Keff}{K_{\mathrm{eff}}}
\newcommand{\kappac}{\kappa_c}
\newcommand{\eps}{\varepsilon}
\newcommand{\betagamma}{\beta\gamma}
\newcommand{\betac}{\beta}
\newcommand{\tauf}{\tau}

% ============================================================
% 定理环境（中文）
% ============================================================
\newtheorem{theorem}{定理}
\newtheorem{lemma}{引理}
\newtheorem{corollary}{推论}
\newtheorem{definition}{定义}
\newtheorem{prediction}{预测}

% ============================================================
% 标题和作者
% ============================================================
\title{\textbf{附录 S：冷原子中光子的负时间延迟作为协调量子动力学的表现：\\ 在 Yakushev 统一协调理论 (YUCT) 框架内的修订分析}}
\author{Alexey V. Yakushev\textsuperscript{1} \\
	\textsuperscript{1}Yakushev Research, \\ 
	\textbf YUCT 理论: \url{https://yuct.org/}, \\ 
	\textbf YPSDC 协议: \url{https://ypsdc.com/}\\
	\textbf DOI: \url{https://doi.org/10.5281/zenodo.18444598}\\
}
\date{2026年5月 \\ \large V.3.0 \\(修正的对数标度与 13\% 温度预测的撤回)}

% ============================================================
% PDF 元数据
% ============================================================
\hypersetup{
	pdftitle={附录 S：冷原子中光子的负时间延迟作为协调量子动力学的表现——在 YUCT 框架内的修订分析},
	pdfauthor={Alexey V. Yakushev},
	pdfsubject={YUCT，负时间延迟，光子，冷原子，量子协调},
	pdfkeywords={YUCT，负时间延迟，光子，冷原子，量子协调，对数标度},
	breaklinks=true,
	pdfencoding=auto,
	bookmarksnumbered=true,
	bookmarksopen=true,
	bookmarksopenlevel=1
}

\begin{document}
	
	\maketitle
	\newpage
	\begin{abstract}
		\noindent
		最新的突破性实验（Angulo 等人，2024）表明，穿越超冷原子云的光子可以以负时间延迟离开介质——波包的峰值比进入时更早出现。这种反直觉的现象虽然不违反因果律，但揭示了深层的量子干涉效应。在本工作中，我们在 Yakushev 统一协调理论（YUCT）框架内重新审视这一解释。基于普适标度律 \(\eps = \kappac \alpha (\ln\Keff)^{\betac}\)，其中 \(\betac = 2/3\) 和 \(\kappac = 1/3\)（推导见附录 X），我们证明了群延迟 \(\tau_g\) 与协调效率 \(\Keff\) 之间的正确关系是对数的：\(|\tau_g| = \alpha_0 \ln\Keff\)。负时间延迟的存在已得到实验确认；已观测到低至 \((-1.14 \pm 0.22)\tau_0\) 的值。然而，与早期的线性假设相反，对数形式意味着 \(\tau_g\) 的温度依赖性极其微弱。使用标准温度退化关系 \(\Keff(T) = K_0 (T/T_0)^{-\gamma}\)，其中 \(\gamma \approx 0.30\)（源自地球物理数据），我们发现温度加倍时 \(|\tau_g|\) 的相对变化量级为 \(10^{-8}\)，远低于当前实验的灵敏度。因此，先前声称的 13\% 效应是错误线性模型的产物，必须予以撤回。我们提供了修正后的理论框架，并讨论了可能对对数标度敏感的其他实验测试，例如改变光学深度或脉冲持续时间。YUCT 的核心预测——负群延迟是协调效率的表现——仍然有效，但其在冷原子系统中的可测试性需要改进的实验设置。
	\end{abstract}
	
	\noindent\textbf{关键词：} YUCT，协调效率，负群延迟，光子-原子相互作用，对数标度，温度无关性，撤回。
	
	\newpage
	\tableofcontents
	\newpage
	
	\section{引言}
	
	光在共振介质中的传播数十年来一直是量子光学的基石。群延迟、脉冲整形和超光速效应在理论和实验上都得到了很好的理解 \cite{Brillouin1960, Milonni2004}。然而，一个基本问题依然存在：光子作为原子激发所花费的时间的物理意义是什么？
	
	Angulo 等人 \cite{Angulo2024} 的一项里程碑式实验通过测量透射光子对独立探测光束引起的交叉克尔相移来探讨这一问题。他们发现，透射光子的积分原子激发时间 \(\tau_T\) 可以为负，即使当 \(\tau_g < 0\) 时也与群延迟 \(\tau_g\) 匹配。具体而言，实验报告了 \(\tau_T/\tau_0 = -1.14 \pm 0.22\)（10 ns 脉冲，光学深度 OD~4）和 \(0.54 \pm 0.28\)（不同参数）的值 \cite{Angulo2024, Nixon2024}。这一非凡的结果挑战了简单的直觉，并要求一个更深层的理论框架。
	
	Yakushev 统一协调理论（YUCT）\cite{Yakushev2026a} 提供了这样的框架。YUCT 假定所有物理相互作用，从量子光学到引力，都是由 \textbf{协调效率} \(\Keff\) 量化的基本协调过程的表现。一个普适标度律将任何相对误差 \(\eps\) 与 \(\Keff\) 联系起来（见附录 X，式~\ref{eq:error-law-corrected}）：
	\begin{equation}
		\eps = \kappac \alpha (\ln\Keff)^{\betac}, \qquad \betac = \frac{2}{3},\ \kappac = \frac{1}{3},
		\label{eq:scaling}
	\end{equation}
	其中 \(\alpha\) 是系统相关的。这一定律已在从 DNA 复制错误到宇宙微波背景涨落的超过 40 个数量级上得到验证 \cite{Yakushev2026c}。在光-物质相互作用的背景下，\(\Keff\) 衡量光子和原子系综之间的相干性。
	
	这里我们表明负时间延迟实验是 YUCT 的直接结果。然而，我们必须修正早期假设 \(\tau_g\) 与 \(\Keff\) 之间为线性关系的尝试。从 YUCT 的信息论结构推导出的正确关系是对数的。这极大地改变了预测的温度依赖性，使其在当前实验精度下可以忽略不计。因此，我们撤回先前发表的 13\% 预测，并提供修订后的自洽分析。
	
	\section{YUCT 的关键概念}
	
	\subsection{协调效率 \(\Keff\)}
	
	在 YUCT 中，任何相互作用系统都由一个字典 \(D\)（共享协议、状态）和一个激活特定协调动作的索引 \(i\) 来表征。协调效率为
	\begin{equation}
		\Keff = \frac{H(\text{动作})}{H(\text{索引})},
	\end{equation}
	其中 \(H\) 表示香农熵。对于量子系统，\(\Keff\) 可以非常大，在完美纠缠的极限下趋于无穷。
	
	\subsection{普适标度律}
	
	涨落和误差服从式~\eqref{eq:scaling}。这一定律源于字典的分形几何，并已在从中子衰变到星系旋转曲线等多样化的系统中得到经验验证 \cite{Yakushev2026c}。对数形式至关重要：它反映了索引空间中的有效分辨率仅随字典大小对数增长的事实。
	
	\subsection{\(\Keff\) 的温度依赖性}
	
	在许多物理系统中，由于热涨落，协调效率随温度升高而降低。从临界现象和分形标度，我们有
	\begin{equation}
		\Keff(T) = K_0 \left(\frac{T}{T_0}\right)^{-\gamma}, \quad \gamma > 0,
		\label{eq:K_T}
	\end{equation}
	其中 \(\gamma\) 是材料相关的指数。乘积 \(\betac\gamma\) 已从地球物理数据中独立确定：\(\betagamma = 0.20 \pm 0.03\) \cite{Yakushev2026b}。这意味着 \(\gamma = \betagamma / \betac = 0.20 / (2/3) = 0.30\)。这个值与无序系统的一般标度论证一致。
	
	\section{光子-原子相互作用的 YUCT 模型}
	
	考虑一个共振光子脉冲入射到冷原子系综。原子充当 \textbf{字典} \(D\)：它们被制备在特定的量子态中（例如通过激光冷却和陷阱）。光子携带一个 \textbf{索引}——其波包形状和频率。相互作用导致一个协调状态，其中光子可能被透射、吸收或散射。
	
	在 Angulo 等人 \cite{Angulo2024} 的实验中，感兴趣的物理量是透射光子的 \textbf{原子激发时间} \(\tau_T\)。根据 YUCT，这个时间与群延迟 \(\tau_g\) 成正比。Thompson 等人 \cite{Thompson2023} 的弱值分析表明，在相关区域中，透射光子的 \(\tau_T = \tau_g\)。
	
	关键步骤是将 \(\tau_g\) 与 \(\Keff\) 联系起来。群延迟源于光子和原子系综之间的相干相互作用。这种相互作用的强度并不与 \(\Keff\) 本身（它可以非常巨大）成正比，而是与光子携带的关于原子状态的 \textbf{信息} 量成正比。在 YUCT 中，相关的信息度量是协调效率的对数，因为索引长度 \(\ell = \log_2 M\) 决定了可区分状态的数量。因此我们假设：
	\begin{equation}
		|\tau_g| = \alpha_0 \, \ln\Keff,
		\label{eq:tau_log}
	\end{equation}
	其中 \(\alpha_0\) 是一个具有时间量纲的基本时间常数。对于铷 D2 线，自然线宽 \(\Gamma = 2\pi \times 6\) MHz 给出特征时间 \(\tau_{\text{nat}} = 1/\Gamma \approx 26\) ns。拟合到实验值 \(|\tau_g| \approx 30\) ns（在 \(K_0 \sim 10^4\) 时）得到 \(\alpha_0 = 30\,\text{ns} / \ln(10^4) = 30 / 9.21 \approx 3.26\) ns。这是合理的：基本相互作用时间在几纳秒的量级。
	
	\subsection{群延迟的温度依赖性}
	
	将式~\eqref{eq:K_T} 代入式~\eqref{eq:tau_log} 得到
	\begin{equation}
		|\tau_g(T)| = \alpha_0 \ln\!\left(K_0 \left(\frac{T}{T_0}\right)^{-\gamma}\right)
		= \alpha_0 \ln K_0 - \alpha_0 \gamma \ln\!\left(\frac{T}{T_0}\right).
	\end{equation}
	记 \(|\tau_g(T_0)| = \alpha_0 \ln K_0\)，我们得到
	\begin{equation}
		|\tau_g(T)| = |\tau_g(T_0)| - \alpha_0 \gamma \ln\!\left(\frac{T}{T_0}\right).
		\label{eq:tau_T_log}
	\end{equation}
	这是核心结果：群延迟的幅度随温度升高而 \textbf{对数} 减小。对于温度比 \(T/T_0 = 2\)，变化为
	\begin{equation}
		\Delta |\tau_g| = \alpha_0 \gamma \ln 2.
	\end{equation}
	使用 \(\alpha_0 \approx 3.26\) ns 和 \(\gamma = 0.30\)，我们得到 \(\Delta |\tau_g| \approx 3.26 \times 0.30 \times 0.693 \approx 0.68\) ns。相对于基线值 \(|\tau_g(T_0)| \approx 30\) ns 的相对变化为 \(0.68 / 30 \approx 2.3\%\)。这仍然在实验的统计不确定度范围内（\(\pm 0.22\) 单位 \(\tau_0\)，即约 \(\pm 5.7\) ns）。然而，该效应处于可探测性的边缘。更重要的是，变化是 \textbf{对数的}，而非幂律，其幅度远小于先前声称的 13\%。
	
	\subsection{与早期（错误的）线性模型的比较}
	
	在早期工作中，假设了线性关系 \(|\tau_g| \propto \Keff\)。这导致 \(|\tau_g| \propto T^{-\gamma}\)，并且当 \(\gamma = 0.30\) 时，温度加倍会产生 13\% 的变化。该预测基于对 \(\Keff\) 在量子系统中作用的误解。正确的对数关系之所以出现，是因为 \(\Keff\) 是熵的比率，而非直接的物理场振幅。线性模型与 YUCT 的信息论基础不相容，在此予以撤回。
	
	\section{基于现有数据的数值一致性检查}
	
	Angulo 等人 \cite{Angulo2024} 的实验提供了一个特定的数据点，使我们能够估计 \(\alpha_0\) 和 \(K_0\)。对于 10 ns 脉冲、光学深度 OD~4 的配置，报告的原子激发时间为 \(\tau_T = -1.14 \tau_0\)。取 \(\tau_0 \approx 26\) ns，我们有 \(|\tau_g| \approx 30\) ns。假设 \(\Keff(T_0) = K_0\) 的量级为 \(10^4\)（典型的冷原子系综有 \(10^6\) 个原子和高相干性），我们得到 \(\alpha_0 = 30\,\text{ns} / \ln(10^4) \approx 3.26\) ns。这个值作为基本光子-原子相互作用的特征时间尺度是合理的。不需要额外的自由参数。
	
	\section{与标准量子光学的比较}
	
	在标准量子光学框架中，共振脉冲的群延迟 \(\tau_g\) 由透射相位对频率的导数给出，\(\tau_g = d\phi/d\omega\)。对于洛伦兹原子响应，这产生一个特征形状，由于反常色散，在共振翼上可能变为负值 \cite{Milonni2004}。然而，标准理论没有预测 \(\tau_g\) 对原子云温度的任何依赖性，除了诸如多普勒展宽（在温度中是线性的，而非对数的）等琐碎效应。观察到的负时间延迟被理解为不涉及能量交换的相干干涉效应。
	
	YUCT 通过将延迟的幅度与协调效率 \(\Keff\) 联系起来，超越了标准理论。对数依赖性是一个独特的特征：如果可以通过改变原子密度、光学深度或脉冲持续时间来改变 \(\Keff\)，延迟应按 \(\ln\Keff\) 标度。这是一个不依赖于温度变化的可检验预测。
	
	\section{修订的实验建议}
	
	由于温度依赖性极弱（对数形式，温度加倍时相对变化仅几个百分点），不太可能用当前精度探测到。因此我们提出其他实验，以更可控的方式直接改变 \(\Keff\)。
	
	\subsection{改变光学深度}
	
	协调效率 \(\Keff\) 与系综中的原子数和相互作用强度成正比。通过减少光学深度（例如降低原子密度），\(\Keff\) 减小。群延迟幅度应遵循 \(|\tau_g| = \alpha_0 \ln\Keff\)。当 \(\Keff\) 减少 10 倍（例如从 \(10^4\) 到 \(10^3\)）时，\(|\tau_g|\) 的变化将为 \(\alpha_0 (\ln 10^4 - \ln 10^3) = \alpha_0 \ln 10 \approx 3.26 \times 2.3026 \approx 7.5\) ns，约为基线值的 25\%，易于测量。这样的实验很简单：对不同原子密度（保持所有其他参数固定）重复测量 \(\tau_T\)，并验证对数标度。
	
	\subsection{改变脉冲持续时间}
	
	光子脉冲的信息内容（索引长度）取决于其时间持续时间和带宽。较短的脉冲携带更多光谱信息，因此具有更大的 \(\Keff\)。通过系统地将脉冲持续时间从 1 ns 改变到 20 ns，可以观察群延迟如何随 \(\ln(\text{带宽})\) 标度。这提供了对数定律的另一个直接检验。
	
	\subsection{温度依赖性作为零检验}
	
	虽然预测的温度效应很小，但并非严格为零。零测量（在实验误差内没有可检测的变化）将与我们的对数模型一致，而早期线性模型会预测明显的 13\% 变化。因此，现有数据已经不利于线性模型，未来的高精度实验可能最终探测到对数偏移。我们鼓励实验者同时检验两者。
	
	\section{与光子质量预测的联系：一致性检查}
	\label{sec:photon_mass_consistency}
	
	YUCT 框架还对光子的静止质量产生了盲目预测，该预测源自同一质量阶梯 \(m = m_e q^{N_f}\)，其中 \(q=(3/2)^{1/3}\)。与实验上限一致的最佳节点是 \(N_\gamma = -410\)，给出
	\[
	m_\gamma \approx 7.3 \times 10^{-19}~\text{eV} \quad (\text{见附录 Photon Mass})。
	\]
	
	人们可能会想知道，这样一个微小的非零光子质量是否会影响冷原子中的群延迟测量。对应于 \(m_\gamma\) 的康普顿波长为
	\[
	\lambda_c = \frac{\hbar}{m_\gamma c} \approx 2.7\times10^{11}~\text{m},
	\]
	它比原子云（典型大小 \(\sim 1\) mm）大许多数量级。因此，由于有限光子质量引起的传播效应在此实验中完全可忽略。负时间延迟纯粹是协调效应，与绝对质量尺度无关。
	
	反过来，一个框架正确预测了以下两者：
	\begin{itemize}
		\item 负群延迟的存在，其幅度由 \(\ln\Keff\) 设定（本附录），
		\item 光子质量的上限（\(m_\gamma < 10^{-18}\) eV）及其离散节点 \(N_\gamma = -410\)（附录 Photon Mass），
	\end{itemize}
	构成了有力的一致性检查。这两个预测涉及截然不同的物理领域（量子光学 vs. 基本粒子性质），却共享相同的底层常数 \(q\)、\(S_{\mathrm{odd}}/S_{\mathrm{even}}\) 以及协调效率 \(\Keff\) 的概念。不需要特设调整来使它们兼容。
	
	因此，负光子延迟实验并不与 YUCT 光子质量预测相矛盾；相反，两者都作为支持统一协调范式的独立支柱。
	
	\section{撤回先前的盲目预测}
	
	在本附录的早期版本（V.2.0）中，我们错误地预测群延迟将按 \(T^{-0.20}\) 标度，导致温度加倍时产生 13\% 的变化。该预测基于错误的线性关系 \(|\tau_g| \propto \Keff\)。这里推导的正确对数关系表明，温度效应大约小两个数量级。因此，关于 13\% 效应的主张是无效的，并正式撤回。我们对此疏忽表示歉意，并感谢数学补丁的共同作者识别出这一不一致性。
	
	然而，YUCT 的核心概念——负群延迟是协调效率的表现——仍然成立。负原子激发时间的实验确认已经支持了光子和原子共享协调字典的观点。修正后的标度定律现在为未来的检验提供了精化的理论框架。
	
	\section{结论}
	
	我们修正了 YUCT 对负光子延迟实验的分析。YUCT 的普适误差定律是对数的，因此群延迟与协调效率之间的关系必须是对数的：\(|\tau_g| = \alpha_0 \ln\Keff\)。这导致非常弱的温度依赖性（对数形式，温度加倍时相对变化仅几个百分点），并消除了先前声称的 13\% 效应。错误的预测被撤回。替代检验，如改变光学深度或脉冲持续时间，提供了对对数标度更灵敏的探测。负时间延迟的存在仍然是量子相互作用协调本质的显著证实，而 YUCT 提供了一致的信息论解释。
	
	\section*{致谢}
	
	作者感谢数学补丁的共同作者指出线性假设与 YUCT 对数结构之间的不一致性。也感谢 YUCT 研讨会参与者的激励性讨论，以及多伦多小组（Aephraim Steinberg、Daniela Angulo 及其同事）的精美实验。
	
	\section*{数据可用性}
	
	所有计算和补充材料可在 \url{https://github.com/Alexey-Yakushev-YUCT/YPSDC} 获取。
	
	\begin{thebibliography}{99}
		
		\bibitem{Angulo2024}
		D. Angulo, K. Thompson, A. Jiao, H. Wiseman, A. Steinberg, ``Experimental evidence that a photon can spend a negative amount of time in an atom cloud,'' arXiv:2409.03680 [quant-ph] (2024).
		
		\bibitem{Nixon2024}
		V.-M. Nixon et al., ``Measuring the time transmitted photons spend as atomic excitations,'' American Physical Society DAMOP Conference, Session Y07 (June 2024).
		
		\bibitem{Yakushev2026a}
		A. V. Yakushev, ``Yakushev Unified Coordination Theory (YUCT),'' Zenodo, \url{https://doi.org/10.5281/zenodo.18444599} (2026).
		
		\bibitem{Yakushev2026b}
		A. V. Yakushev, ``Appendix P: Temperature Dependence of Gravity,'' Zenodo (2026).
		
		\bibitem{Yakushev2026c}
		A. V. Yakushev, ``Appendix L: Fractal Error Scaling,'' Zenodo (2026).
		
		\bibitem{Yakushev2026d}
		A. V. Yakushev, ``Appendix X: Generalization of Shannon Information Theory,'' Zenodo (2026).
		
		\bibitem{Thompson2023}
		K. Thompson et al., ``How much time does a photon spend as an atomic excitation before being transmitted?'' arXiv:2310.00432 (2023).
		
		\bibitem{Brillouin1960}
		L. Brillouin, \textit{Wave Propagation and Group Velocity} (Academic Press, 1960).
		
		\bibitem{Milonni2004}
		P. Milonni, \textit{Fast Light, Slow Light and Left-Handed Light} (CRC Press, 2004).
		
		\bibitem{Wiseman2023}
		H. M. Wiseman, A. M. Steinberg, and M. Hallaji, ``Obtaining a single-photon weak value from experiments using a strong coherent state,'' AVS Quantum Sci. \textbf{5}, 024401 (2023).
		
	\end{thebibliography}
	
\end{document}